% ============================================================
% CAPÍTULO 5
% ============================================================
\capitulo{Resultados y Discusión}

\section{Rendimiento del Modelo}
Sobre 359 clips de 133 caballos, el modelo MIL multimodal obtuvo un \fmacro{} de 0,740 ($DE = 0{,}052$ entre particiones; IC 95~\% por \textit{bootstrap} de caballos [0,677; 0,797]) y una exactitud de 0,744. Supera el criterio de aceptación fijado antes de entrenar (0,691). La Tabla~\ref{tab:porclase} detalla las métricas por clase y la Tabla~\ref{tab:confusion}, la matriz de confusión.

\begin{table}[H]
\caption{Métricas por clase del modelo MIL multimodal}\label{tab:porclase}
{\singlespacing
\begin{tabular}{@{}lrrrr@{}}
\toprule
Clase & Precisión & Sensibilidad & \textit{F}\textsubscript{1} & \textit{n} \\
\midrule
\clase{Disgustado} & 0,740 & 0,689 & 0,714 & 103 \\
\clase{Interesado} & 0,764 & 0,797 & 0,780 & 118 \\
\clase{Sereno} & 0,729 & 0,739 & 0,734 & 138 \\
\bottomrule
\end{tabular}}
\notatabla{Predicciones fuera de muestra de las cuatro particiones agrupadas por caballo.}
\end{table}

\begin{table}[H]
\caption{Matriz de confusión del modelo MIL multimodal}\label{tab:confusion}
{\singlespacing
\begin{tabular}{@{}lrrr@{}}
\toprule
& \multicolumn{3}{c}{Clase predicha} \\
\cmidrule(l){2-4}
Clase real & \clase{Disgustado} & \clase{Interesado} & \clase{Sereno} \\
\midrule
\clase{Disgustado} & 71 & 13 & 19 \\
\clase{Interesado} & 5 & 94 & 19 \\
\clase{Sereno} & 20 & 16 & 102 \\
\bottomrule
\end{tabular}}
\end{table}

El error dominante es la confusión entre \clase{Sereno} y \clase{Disgustado} (20 y 19 clips), más frecuente que la que parecería semánticamente más cercana. Algunos caballos son clasificados mal en todos sus clips; un error sistemático sobre un mismo animal es más compatible con una etiqueta discutible o con una apariencia atípica que con ruido del modelo, pero sin una segunda opinión sobre las etiquetas no es posible distinguir entre ambas causas.

\section{Variabilidad entre Particiones}
El \fmacro{} de prueba fue de 0,718, 0,769, 0,677 y 0,795 en las cuatro particiones. La dispersión indica que el rendimiento depende de qué caballos quedan en prueba, lo esperable cuando el tamaño efectivo de la muestra son 133 animales y no 359 clips independientes.

\section{Líneas de Base}
La Tabla~\ref{tab:lineas} compara el modelo con las líneas de base evaluadas sobre las mismas particiones.

\begin{table}[H]
\caption{Comparación con líneas de base}\label{tab:lineas}
{\singlespacing
\begin{threeparttable}
\begin{tabular}{@{}lrr@{}}
\toprule
Modelo & \fmacro{} (\textit{DE}) & Exactitud \\
\midrule
Clase mayoritaria & 0,183 (0,033) & 0,384 \\
Azar estratificado\tnote{a} & 0,325 (0,025) & --- \\
Duración del clip & 0,544 (0,035) & 0,532 \\
Geometría sola (media y desvío) & 0,607 (0,033) & 0,618 \\
Metadatos de filmación & 0,618 (0,034) & 0,635 \\
Visual, un fotograma central & 0,679 (0,024) & 0,688 \\
Visual, promedio de 16 fotogramas & 0,741 (0,020) & 0,749 \\
MIL multimodal (modelo final) & 0,740 (0,052) & 0,744 \\
\bottomrule
\end{tabular}
\begin{tablenotes}[flushleft]\footnotesize
\item[a] Esperanza del \fmacro{} de predicciones aleatorias con las proporciones de clase; percentil 95: 0,367.
\end{tablenotes}
\end{threeparttable}}
\notatabla{Todas las líneas de base son regresiones logísticas evaluadas con las mismas cuatro particiones por caballo que el modelo final. \textit{DE} = desvío estándar entre particiones.}
\end{table}

Tres resultados se desprenden de la comparación. Primero, el modelo supera con amplitud al azar y a la clase mayoritaria. Segundo, los metadatos de filmación por sí solos alcanzan 0,618: la forma en que se filmó cada clase contiene señal. El modelo los supera por 0,121 (IC 95~\% [0,031; 0,211]), de modo que su rendimiento no se reduce a esos metadatos. Tercero, y más importante, un clasificador lineal sobre el \textit{promedio} de las características visuales de los 16 fotogramas obtiene 0,741, prácticamente lo mismo que el modelo final: la diferencia es de $-0{,}005$ con IC 95~\% [$-0{,}068$; 0,058], que incluye el cero. Usar 16 fotogramas sí aporta respecto de uno solo (diferencia de 0,059, IC 95~\% [0,001; 0,116]), pero no se demuestra que la atención, la geometría ni el orden agreguen algo sobre un resumen estático de esos fotogramas.

\section{Control del Contexto de Filmación}
Para estimar cuánto del rendimiento depende de rasgos de la filmación, se restringió la evaluación a los clips que comparten resolución ($1080 \times 1920$) y tienen al menos 5~s de duración, sin reentrenar. Sobre esos 236 clips el modelo conserva un \fmacro{} de 0,705 y una exactitud de 0,746. El subconjunto contiene sólo 28 clips de \clase{Disgustado}, por lo que el resultado es indicativo; aun así, la caída es moderada y sugiere que el modelo no depende principalmente de la duración ni de la resolución.

Este control no alcanza al contexto que la etiquetadora describe (box de hipódromo para \clase{Interesado}, interacción en manada para \clase{Disgustado}). Que la clase se asocie al entorno es inherente al fenómeno, pero el modelo puede aprovecharlo como pista: los videos con varios animales en cuadro pertenecen casi siempre a \clase{Disgustado} y el modelo los clasifica mejor que al resto de esa clase.

\section{Sobreajuste}
El registro por época del entrenamiento interno muestra un sobreajuste marcado: el \fmacro{} medio fue de 0,979 en entrenamiento, 0,712 en validación y 0,740 en prueba (Anexo~C). La parada temprana contiene el problema: el rendimiento de prueba se mantiene en línea con el de validación. Con 132\,355 parámetros y unos 200 clips de entrenamiento por partición, la memorización del conjunto de entrenamiento es esperable.

\section{Evaluación de la Hipótesis Temporal}
La hipótesis temporal no queda respaldada. La GRU y la atención con codificación posicional, sensibles al orden, no superaron de manera consistente a las variantes invariantes. En el modelo de promedio, permutar los fotogramas no cambia el resultado por construcción, y el modelo final, cuya atención no usa la posición, es igualmente invariante al orden. La comparación con el promedio estático refuerza la conclusión: un resumen sin orden ni atención alcanza el mismo rendimiento.

La conclusión es acotada. No implica que la dinámica sea irrelevante en la expresión equina ---en dolor equino se reportó lo contrario \parencite{broome2019}---, sino que \textit{con este conjunto, esta ventana de hasta 7~s y estas representaciones, no se encontró evidencia de que el orden de los 16 fotogramas mejore la clasificación}.

\section{Variantes No Adoptadas}
El ajuste fino parcial de la ResNet18 (última etapa convolucional) obtuvo 0,747 ($DE = 0{,}057$), una diferencia pequeña frente a la variabilidad entre particiones que no justificaba el costo ni el riesgo adicional de sobreajuste. Tampoco mejoraron el modelo un canal geométrico con más puntos anatómicos, puntos de boca y ollares, la validación cruzada anidada para elegir la época, el enmascarado de clips incompletos ni una representación del movimiento de la cabeza. La geometría resultó útil sólo como complemento: el MIL visual sin geometría obtuvo 0,686.

\section{Limitaciones}
\begin{enumerate}
  \item \textbf{Un único observador.} El acuerdo se midió en un piloto sobre imágenes, no sobre este conjunto. Se desconoce el techo de rendimiento que impone la ambigüedad de las etiquetas.
  \item \textbf{El criterio de etiquetado no es sólo facial.} Incluye ambiente, comportamiento y postura, información a la que el modelo, que ve un recorte de la cara, no accede por completo.
  \item \textbf{Asociación entre clase y contexto.} La clase se relaciona parcialmente con el establecimiento, la situación y la duración de la toma.
  \item \textbf{Tamaño efectivo.} Los 133 caballos limitan la precisión de las estimaciones, como muestra la variabilidad entre particiones.
  \item \textbf{Sin validación externa.} No se evaluó el modelo con videos de otras fuentes, cámaras o regiones.
  \item \textbf{Validación interna pequeña.} La elección de la época se apoya en pocos caballos por partición.
  \item \textbf{Pose sin adaptación.} DeepLabCut se usa sin entrenamiento adicional sobre este conjunto.
  \item \textbf{Costo de inferencia.} Procesar un video nuevo demora de uno a tres minutos con GPU.
\end{enumerate}

\section{Discusión General}
Los resultados muestran que la cara del caballo, tal como la captura un video corto de celular, contiene señal suficiente para discriminar las tres clases definidas en el conjunto, y que esa señal se generaliza a animales no vistos durante el entrenamiento. El rendimiento es comparable entre clases y supera claramente a las líneas de base que no miran la cara.

Al mismo tiempo, el trabajo muestra que la complejidad del modelo no fue lo que determinó el resultado. Los avances más grandes provinieron de decisiones metodológicas: eliminar la fuga por identidad, corregir la orientación de los videos, ampliar el encuadre para incluir las orejas, rediseñar la selección de fotogramas y unificar entrenamiento e inferencia. Un clasificador lineal sobre un resumen estático de varios fotogramas iguala a la arquitectura final. Con un conjunto de este tamaño, la calidad de los datos y el control de sesgos pesan más que la arquitectura.
